\begin{definition}[Measure root domain-ready package]
\label{def:measure-root-domain-ready-package}
The $\MeasureUp$ root domain-ready package is the paper-side root package
whose fields are the following ten rows:
\begin{enumerate}
  \item the traditional constructive measure-space target, namely a
  measurable-event carrier, empty event, finite unions, countable unions,
  relative differences, disjointness witnesses, nonnegative $\RealUp$
  endpoints, and countable-sum endpoints;
  \item the BEDC source row over $\mathsf{BHist}$, together with
  $\NameCert_{\RealUp}$ endpoint comparison, the additive $\AbGroupUp$ real
  row, and $\StreamNameUp$ countable-family names;
  \item the carrier rows displayed by
  \autoref{def:measure-sigma-algebra-carrier-surface};
  \item the componentwise classifier rows displayed by
  \autoref{def:measure-certificate-surface};
  \item the stability rows for measurable events, endpoint comparisons,
  countable-family presentations, relative differences, and countable sums;
  \item the visible ledger and exactness rows for hidden membership,
  disjointness, countable-union witnesses, and endpoint readback;
  \item the core theorem list consisting of carrier-classifier transport,
  sigma-additivity, empty-zero, finite disjoint-union readback,
  finite-support countable readback, relative-difference additivity, and
  continuity from below;
  \item the exactness predicate stating that every consumer row is one of the
  displayed root rows;
  \item the explicit absence of any $\StdBridge$ row;
  \item the status row exposing this object as the paper-side root package
  for $\MeasureUp$.
\end{enumerate}
\end{definition}

\begin{theorem}[Measure root domain carrier-classifier ready]
\label{thm:measure-root-domain-carrier-classifier-ready}
The domain-ready package of
\autoref{def:measure-root-domain-ready-package} supplies carrier transport,
componentwise classifier transport, visible event-row transport, ledger
transport, and $\RealUp$ endpoint transport for the root $\MeasureUp$ domain.
These rows are exactly the carrier-classifier part of the root public surface
under the input scope of $\mathsf{BHist}$, $\hsame$, $\Cont$,
$\NameCert_{\RealUp}$, the additive $\AbGroupUp$ real row, and
$\StreamNameUp$ countable-family names.
\end{theorem}
\begin{proof}
Start from the carrier and classifier theorem
\autoref{thm:measure-root-carrier-classifier-public-obligation}. It projects
the carried $\MeasureUp$ code, visible event rows, ledger rows, and $\RealUp$
endpoint rows along the public classifier.

Apply \autoref{thm:measureup-classifier-stability-obligations}. Its
certificate-level projections identify the same componentwise classifier
movement for event presentations, values, and countable-sum components.

Finally apply \autoref{thm:measure-root-input-dependency-scope}. The input
scope pins the transport rows to $\mathsf{BHist}$, $\hsame$, $\Cont$,
$\NameCert_{\RealUp}$, the additive real row, and the endpoint-family
$\StreamNameUp$ rows, with no downstream consumer row used as a hypothesis.
These projections are precisely the carrier-classifier fields of the
domain-ready package.
\end{proof}

\begin{theorem}[Measure root domain sigma-additive ready]
\label{thm:measure-root-domain-sigma-additive-ready}
The domain-ready package of
\autoref{def:measure-root-domain-ready-package} supplies sigma-additivity,
empty-zero, finite disjoint-union readback, finite-support countable
readback, relative-difference additivity, and continuity-from-below rows for
the root $\MeasureUp$ domain.
\end{theorem}
\begin{proof}
Apply \autoref{thm:measure-root-sigma-additive-public-obligation}. It gives
the public sigma-additive comparison, the empty-zero row, finite
disjoint-union readback, and the relative-difference row from the root
NameCert surface.

Next use \autoref{thm:measure-finite-support-countable-readback}. This reads
finite-support countable families through the same countable-union,
disjointness, endpoint, and additive rows.

Apply \autoref{thm:measure-relative-difference-additivity} to keep the
relative-difference row visible as a root domain consequence, and apply
\autoref{thm:measure-continuity-from-below-integral-scope} to expose
continuity from below at the same measure surface.

The listed rows use the carrier, endpoint, countable-family, additive, and
ledger fields already named by the domain-ready package, so they repack as
its sigma-additive face.
\end{proof}

\begin{theorem}[Measure root domain downstream ready]
\label{thm:measure-root-domain-downstream-ready}
The carrier-classifier theorem
\autoref{thm:measure-root-domain-carrier-classifier-ready} and the
sigma-additive theorem
\autoref{thm:measure-root-domain-sigma-additive-ready} expose exactly the
rows consumed by $\IntegralUp$, $\ProbSpaceUp$, and
$\FunctionalAnalysisUp$. The resulting downstream face contains no quotient
event equality, no host set equality, and no standard bridge row.
\end{theorem}
\begin{proof}
First apply \autoref{thm:measure-root-downstream-unblock-package}. It reads
the carrier-classifier rows and sigma-additive rows as the public
$\MeasureUp$ package available to the three downstream consumers.

For $\IntegralUp$, apply \autoref{thm:measure-integral-consumer-rows}. The
consumed rows are the covered measurable-event rows, $\RealUp$ endpoint
classifier rows, countable-sum endpoint rows, head-tail readback rows, and
finite disjoint-union readback rows.

For $\ProbSpaceUp$, apply \autoref{thm:measure-probability-consumer-rows}.
The consumed rows are total-event normalization together with empty-zero and
finite disjoint-union readback.

For $\FunctionalAnalysisUp$, use the consumer projection recorded in
\autoref{def:measure-root-consumer-package} and made explicit by
\autoref{thm:measure-functionalanalysis-consumer-rows}. Its rows are the
nonnegative endpoint, monotone inclusion, relative-difference, and
countable-sum classifier rows.

The exactness predicate in \autoref{def:measure-root-consumer-package}
excludes every row not named above. Hence the downstream face of the
domain-ready package contains neither quotient event equality nor host set
equality, and its standard-bridge field remains the explicit no-$\StdBridge$
row.
\end{proof}
